\section{Reverse a Linked List (Recursive)}
\label{sec:ll-reverse-recursive}

\problemstatement{Reverse a singly linked list using recursion, returning
the new head.}

\begin{patternbox}
The recursive reversal expresses the same rewiring differently: reverse the
rest of the list first, then fix the link between the current node and its
successor on the way back up. The new head bubbles up unchanged from the
deepest call.
\end{patternbox}

\subsection*{Reverse the tail, then flip one link}

Recurse to the end: the base case (empty or single node) returns itself as
the new head. On the way back, the current node's successor
(\texttt{head->next}) should now point back at the current node, so set
\texttt{head->next->next = head} and \texttt{head->next = null}. The deepest
call's return value -- the original last node -- is the new head and is
passed up unchanged through every frame.

\begin{ideabox}[Fix the Link on the Return Trip]
Going down, we only recurse; coming up, each frame flips exactly one link
between \texttt{head} and \texttt{head->next}. Because the new head is
determined at the very bottom, it is threaded back up untouched while the
per-frame flips reverse the pointers.
\end{ideabox}

\subsection*{Algorithm}

\begin{enumerate}
  \item If \texttt{head} is null or \texttt{head->next} is null, return
        \texttt{head}.
  \item \texttt{newHead = reverse(head->next)}.
  \item \texttt{head->next->next = head}; \texttt{head->next = null}.
  \item Return \texttt{newHead}.
\end{enumerate}

\subsection*{Dry run}

\dryrun{List $1\to2\to3$.}

\begin{itemize}
  \item Recurse to $3$ (base) $\to$ newHead $=3$.
  \item At $2$: $3\to2$, $2\to$null. At $1$: $2\to1$, $1\to$null.
  \item Result $3\to2\to1$.
\end{itemize}

\subsection*{C++ implementation}

\begin{lstlisting}
#include <bits/stdc++.h>
using namespace std;

struct ListNode {
    int val; ListNode* next;
    ListNode(int v) : val(v), next(nullptr) {}
};

ListNode* build(const vector<int>& a) {
    ListNode dummy(0), *t = &dummy;
    for (int x : a) { t->next = new ListNode(x); t = t->next; }
    return dummy.next;
}

ListNode* reverse(ListNode* head) {
    if (!head || !head->next) return head;        // base case
    ListNode* newHead = reverse(head->next);
    head->next->next = head;                      // flip one link
    head->next = nullptr;
    return newHead;                               // threaded up unchanged
}

void print(ListNode* h) { for (; h; h = h->next) cout << h->val << " "; cout << "\n"; }

int main() {
    ListNode* head = build({1, 2, 3, 4});
    head = reverse(head);
    print(head);                                  // 4 3 2 1
    return 0;
}
\end{lstlisting}

\begin{complexitybox}
\textbf{Time Complexity.} $O(n)$. \textbf{Space Complexity.} $O(n)$ for the
recursion stack (versus $O(1)$ for the iterative form).
\end{complexitybox}

\begin{takeawaybox}
Recursive reversal reverses the tail then flips one link per return frame,
carrying the new head up from the base case. It is more elegant but uses
$O(n)$ stack; prefer the iterative version when depth could overflow.
\end{takeawaybox}
